\documentclass[11pt,a4paper]{article}
\usepackage[margin=2.5cm]{geometry}
\usepackage{enumitem}
\usepackage{booktabs}
\usepackage{hyperref}
\usepackage{titlesec}
\usepackage{fancyhdr}
\usepackage{xcolor}
\usepackage{tabularx}
\usepackage{longtable}

\hypersetup{
    colorlinks=true,
    linkcolor=blue!60!black,
    urlcolor=blue!60!black
}

\titleformat{\section}{\large\bfseries}{}{0em}{}
\titleformat{\subsection}{\normalsize\bfseries}{}{0em}{}

\pagestyle{fancy}
\fancyhf{}
\rhead{\small\thepage}
\lhead{\small MA Asian Studies --- Thesis Protocol}
\renewcommand{\headrulewidth}{0.4pt}

\setlength{\parindent}{0pt}
\setlength{\parskip}{0.8em}

\begin{document}

\begin{center}
{\Large\bfseries MA THESIS PROTOCOL 2025--2026}\\[0.3em]
{\Large\bfseries Asian Studies Program}\\[0.3em]
{\normalsize Leiden University $\cdot$ Faculty of Humanities}\\[1.5em]
\end{center}

\begin{quote}
\small\textit{This document is adapted from the MA Asian Studies thesis protocol for student reference. Some content intended for lecturers and supervisors has been removed. For the most current information, always consult your supervisor and check the \href{https://studiegids.universiteitleiden.nl/en/courses/134994/ma-thesis-asian-studies-60-ec}{MAAS e-guide}.}
\end{quote}

\bigskip
\hrule
\bigskip

\section{Assignment of a Supervisor}

New students in the MA60 and MA120 are required to submit a thesis distribution form shortly after the start of their program. The form includes the names of potential supervisors and a description of your research interests. It is due on \textbf{September 24th} (or early February if you started in February).

By mid-October (end of February for spring entry), the chair of the program matches students with a supervisor.

\section{Thesis and Methodology Sessions}

As part of the Introduction to Asian Studies course, you will receive three lectures covering methodology and thesis proposal preparation. Make sure to complete the \textit{assignment} on Brightspace before the first session. Students joining in the spring will be offered these sessions in the first three weeks of the semester and will have to complete Introduction to Asian Studies in the fall.

\section{Thesis Supervision and Proposal}

You should schedule a first meeting with your supervisor in November (or March if you joined in the spring), soon after the appointment of the thesis supervisor. Use this meeting to discuss your interests and plans and to receive guidance on sources, the state of the field, and suitable research questions.

Following this meeting, you prepare a \textbf{thesis proposal} which is due on \textbf{January 5th} (or \textbf{June 1st} if you started in semester 2). The proposal should be:

\begin{itemize}
    \item Approximately \textbf{1,500 words} including a bibliography
    \item State a clear research question and describe the relevant literature and methodology
    \item Submitted by email to the supervisor
\end{itemize}

Your supervisor will provide written feedback by \textbf{January 20th} (June 20th in semester 2) to allow sufficient time for revisions.

After receiving feedback, the final thesis proposal must be submitted to the supervisor by email by \textbf{January 31st} (for February intake: \textbf{July 1st}).

\textbf{The proposal is graded on a pass/fail basis.} Students who fail may submit a revised proposal. For students who submitted their final proposal on January 31st, resubmission is possible from April 1st.

The thesis proposal assessment is a required component of the course Introduction to Asian Studies. Failure to complete the thesis proposal within the academic year means students must repeat the course.

\section{Writing the Thesis and Draft Submission}

After the approval of the thesis proposal, you will start writing draft chapters in consultation with your supervisor and meet with your supervisor at least twice.

There are two deadlines for the MA thesis:

\textbf{First submission deadline:} \textbf{May 15th} (or \textbf{November 1st} for the February intake). Students who submit their thesis, or part of it, to their supervisor on or before the first deadline \textbf{have a right to receive feedback}. They will receive their papers back, with commentary, within 10 working days. This is the final round of feedback provided by supervisors.

Students who do not submit by the first deadline forfeit the right for further feedback. They retain the right to submit a final version by the final deadline.

\textbf{Final submission deadline:} The final version of the thesis must be submitted by email to the supervisor, second reader, and to \href{mailto:MAthesis@hum.leidenuniv.nl}{MAthesis@hum.leidenuniv.nl} on or before \textbf{July 1st} (December 15th for spring entry).

If you face extenuating circumstances or need an extension, you must submit your request via the examinations committee. Extensions are not individually granted by supervisors.

\section{Summary of Student Deadlines}

\begin{center}
\begin{tabularx}{\textwidth}{Xll}
\toprule
\textbf{Milestone} & \textbf{Fall entry} & \textbf{Spring entry} \\
\midrule
Thesis distribution form due & September 24 & Mid-February \\
First meeting with supervisor & November & March \\
Thesis proposal draft due & January 5 & June 1 \\
Thesis proposal final draft due & January 31 & July 1 \\
Thesis draft due (first submission) & May 15 & November 1 \\
\textbf{Thesis final draft due} & \textbf{July 1} & \textbf{December 15} \\
\midrule
Resit thesis draft & November 1 & May 15 \\
\textbf{Resit thesis final draft} & \textbf{December 15} & \textbf{July 1} \\
\bottomrule
\end{tabularx}
\end{center}

\section{Thesis Second Attempt}

Students who do not pass the first assessment, or who fail to complete the thesis by the final deadline, may submit a new draft by the resit draft deadline and a final version by the resit final deadline, under the guidance of the same supervisor.

\section{Thesis Requirements and Format}

The final version should be \textbf{between 12,000 and 15,000 words}, including footnotes and references, but excluding appendices.

You should use an established reference format according to the standards of the discipline (see, for example, \textit{Journal of Asian Studies} for social science or \textit{The American Historical Review} for history). Use that format \textbf{consistently and correctly} throughout the thesis.

There is no single standard format for the layout (font, font size, margins, spacing, footnotes or endnotes, etc.), so please consult your supervisor about their preference.

\section{Assessment Criteria}

Your thesis is assessed by your supervisor (first reader) and a second reader. The assessment is based on the following criteria:

\subsection{Knowledge and Insight}

The research question is based on a problem that reflects insight into the key discussions and methods of the field. Assessment also considers whether the problem is clear, relevant, well defined, placed in the existing literature, and original.

\subsection{Application of Knowledge and Insight}

The thesis critically analyzes primary material and sources. It should put complex concepts into practice, use effective research methods, and draw on secondary sources meant for an advanced academic audience. It should also describe and justify the adopted method and apply knowledge and insight in broader or multidisciplinary contexts.

\subsection{Reaching Conclusions}

Assessment considers logical and consistent reasoning, well-founded conclusions, and the degree to which the thesis question is answered. It also considers connections to other and future research, attention to social and ethical aspects, and critical reflection on one's role as researcher.

\subsection{Communication}

Assessment considers language use, including linguistic competence, readability, style, spelling, grammar, and the use and explanation of correct terminology. It also considers the structure and layout of the thesis, as well as the apparatus, including annotations, correct use of reference guidelines, completeness of references, and bibliography.

\subsection{Learning Skills (Process)}

Assessment considers the degree of independence, planning and time management, handling of supervisor feedback, and, if applicable, participation in a thesis group.

\section{Specialization Requirements}

In addition to the general criteria, your thesis must meet the requirements of your specialization.

\textbf{MA60 East Asian Studies:} Broad knowledge of East Asia. Knowledge of the history, development, and current trends in East Asian Studies. The ability to locate, assess, and use original sources in modern or classical Chinese, modern or classical Japanese, or modern or classical Korean.

\textbf{MA60 South Asian Studies:} Broad knowledge of South Asia. Knowledge of the history, development, and current trends in South Asian Studies. For students who have completed intermediate or advanced language electives, the ability to use original sources in Hindi, Sanskrit, or Classical Tibetan.

\textbf{MA60 Southeast Asian Studies:} Broad knowledge of Southeast Asia. Knowledge of the history, development, and current trends in Southeast Asian Studies. For students who have completed intermediate or advanced language electives, the ability to use original sources in Malay/Indonesian.

\textbf{MA60 History, Arts and Culture:} Broad knowledge of one or more regions in Asia. Knowledge of history, development, and current trends in the study of the history, art, or culture of these regions. Familiarity with disciplinary theories and methods associated with the study of history, arts, and/or culture.

\textbf{MA60 Politics, Society, and Economy:} Broad knowledge of one or more regions in Asia. Knowledge of history, development, and current trends in the study of the politics, society, or economy of these regions. Familiarity with disciplinary theories and methods associated with the study of politics, society, and/or economy.

\textbf{MA120 Chinese, Japanese, Korean Studies:} Foundational knowledge of the history and development of Asian Studies. Knowledge of research in social sciences and humanities disciplines grounded in Asian Studies. Understanding of the concepts, terminology, and approaches relevant to the regional and disciplinary specialization. Knowledge of the main research methods and techniques and their application. Skills in locating, retrieving, assessing, and using sources and specialist secondary literature for original research. The ability to independently design, conduct, and complete original research. Excellent command of Chinese, Japanese, or Korean that enables use of the language for academic and professional purposes.

\end{document}
